\section*{Chapter \ref{ch:b3:binary-outcomes} --- Binary Outcomes and Maximum Likelihood}

\begin{solution}{exo:b3:binary-outcomes:1}
They differ most at the edges: at zero earnings the linear model predicts 0.704,
the logit 0.665 and the probit 0.681; near the top the linear model reaches 1
while the curves flatten at about 0.97. % figdata: work_fits
All three fail for the men with no earnings in 1975, whose observed share is
0.438. % figdata: work_bins
Those men form a group of their own --- most were out of work --- and a smooth
function of earnings that also fits the men with earnings cannot drop to them.
\end{solution}

\begin{solution}{exo:b3:binary-outcomes:2}
$\Lambda(0.8) = 1/(1 + \eu^{-0.8}) = 0.690$; marginal effect $0.690 \times 0.310
\times 0.5 = 0.107$. $\Lambda(3) = 0.953$; marginal effect $0.953 \times 0.047
\times 0.5 = 0.022$ (0.023 from the unrounded values). The same coefficient moves the probability a fifth as much
near the top.
\end{solution}

\begin{solution}{exo:b3:binary-outcomes:3}
$\eu^{-0.2981} = 0.742$: other things equal, a married man's odds of having
earnings in 1978 are 0.74 times an unmarried man's, about a quarter lower.
% output: lpm_logit_probit
\end{solution}

\begin{solution}{exo:b3:binary-outcomes:4}
IIA says the tram takes the same percentage of each mode's users. Data showing
that the bus lost a large percentage of its riders while driving lost a small
one (say the bus share fell by 40~\% and the car share by 5~\%) contradict it, as
does a change in the ratio of car to bus riders after the tram opened.
\end{solution}

\begin{solution}{exo:b3:binary-outcomes:5}
The tested solution is \code{ex05_zero_earnings.py}, in the chapter's
\code{exercises/} folder.
\begin{enumerate}
\item The \omterm{def:b3:binary-outcomes:mle}{log-likelihood} rises from $-5{,}507.9$ to $-5{,}149.5$, and the mean
fitted probability for men with no earnings in 1975 falls from 0.665 to 0.438,
exactly the observed share. % output: ex05_zero_earnings
\item The \omterm{def:b3:binary-outcomes:ame}{average marginal effect} of the dummy is $-0.135$: having no earnings in
1975 lowers the probability of earnings in 1978 by 13.5 points beyond what a
low income predicts. That of a thousand dollars of 1975 earnings falls to
0.0082.
\end{enumerate}
\end{solution}

\begin{solution}{exo:b3:binary-outcomes:6}
$\ln 0.8 + \ln 0.7 + \ln 0.6 + \ln 0.9 + \ln 0.5 = -0.223 - 0.357 - 0.511 - 0.105
- 0.693 = -1.889$. (For $y = 0$ the term is $\ln(1 - p)$.) The fifth observation
contributes most, $-0.693$: its outcome had a fitted probability of only one
half.
\end{solution}

\begin{solution}{exo:b3:binary-outcomes:7}
A 1~\% rise in product 1's price raises the shares of products 2 and 3 by 0.21~\%
each; a 1~\% rise in product 4's price raises them by 0.47~\% each.
% output: logit_demand
The cross elasticity with respect to product $k$'s price is $-\alpha p_k s_k$, which
does not depend on the product whose share responds: IIA.
\end{solution}

\begin{solution}{exo:b3:binary-outcomes:8}
\begin{enumerate}
\item $\ell(b) = n_1 b - n\ln(1 + \eu^b)$; the score is $n_1 - np$ and the
second derivative $-np(1 - p)$, with $p = \Lambda(b)$. The Newton step is $b +
(n_1 - np)/(np(1 - p))$.
\item From $b = 0$: $p = 0.5$, $b_1 = 0 + 30/25 = 1.2$. Then $p = \Lambda(1.2) =
0.7685$ and $b_2 = 1.2 + (80 - 76.85)/(100 \times 0.7685 \times 0.2315) = 1.2 +
0.177 = 1.377$, against $\ln 4 = 1.386$.
\end{enumerate}
\end{solution}

\begin{solution}{exo:b3:binary-outcomes:9}
First, sort the sample by fitted value, cut it into ten groups and compare each
group's mean outcome with its mean fitted value: a linear model that is wrong
will be off systematically at the ends. Second, count fitted values outside
$[0, 1]$: 15.7~\% exceed one here. % output: lpm_logit_probit
A third check adds the squared fitted values as a regressor; a significant
coefficient signals curvature the model misses.
\end{solution}

\begin{solution}{exo:b3:binary-outcomes:10}
The score equation for the constant is $\sum_i (y_i - p_i) = 0$, so $\frac1n\sum_i
p_i = \frac1n\sum_i y_i$. With only group dummies and no constant, every member
of group $g$ has the same index $\beta_g$ and probability $p_g = \Lambda(\beta_g)$;
the score equation for $\beta_g$ is $\sum_{i \in g} (y_i - p_g) = 0$, so $p_g$ is
the group's share. In \cref{exo:b3:binary-outcomes:5} the dummy for no earnings
plays this role, together with the constant: the fitted share for that group
equals the observed 0.438 whatever the other regressors.
\end{solution}

\begin{solution}{exo:b3:binary-outcomes:11}
The tested solution is \code{ex11_endogenous_price.py}, in the chapter's
\code{exercises/} folder.
\begin{enumerate}
\item OLS gives $-1.071$ against a true $-2$. % output: ex11_endogenous_price
It is biased toward zero: products with high unseen quality have both high
prices and high shares, so price seems to hurt sales less than it does.
\item \omterm{def:b3:instrumental-variables:tsls}{Two-stage least squares} with the cost as instrument gives $-2.033$
(s.e.\ 0.064).
\end{enumerate}
\end{solution}

\begin{solution}{exo:b3:binary-outcomes:12}
\begin{enumerate}
\item Each existing share is multiplied by 0.9: car 54~\%, bus 27~\%, cycling
9~\%, cycle lane 10~\%.
\item The lane is a close substitute for cycling and a poor one for driving. It
should take mostly from cyclists and some bus riders: IIA overstates the fall in
driving and understates the fall in cycling on the old roads.
\end{enumerate}
\end{solution}

\begin{solution}{pb:b3:binary-outcomes:bus}
\textbf{1.} $\Pr(\text{car}) = \Lambda(0.6131 - 0.3198\,\Delta\text{cost} -
0.0537\,\Delta\text{time})$, with differences car minus bus.
% figdata: commuters_binary

\textbf{2.} $\Delta\text{cost} = 2$, $\Delta\text{time} = -15$: $x'\hat\beta =
0.6131 - 0.6396 + 0.8055 = 0.779$; $\Pr = \Lambda(0.779) = 0.685$.

\textbf{3.} At equal cost and time, the log odds of driving are 0.61: a
probability of $\Lambda(0.6131) = 0.65$. It measures what the model leaves out
(comfort, flexibility) relative to the bus; a different base would change its
meaning but not the probabilities.

\textbf{4.} $t = -0.3198/0.0546 = -5.86$: clearly different from zero.

\textbf{5.} Only the ratio of the coefficients to the spread of the taste shocks
is identified. The logit fixes the spread (variance $\pi^2/3$) so that the
coefficients are defined; ratios such as the value of time do not depend on
the choice.

\textbf{6.} If the three-way logit is right, IIA implies that among those who
chose the car or the bus, the choice between them follows a binary logit with
the same coefficients; dropping the cyclists loses information but not
consistency.

\textbf{7.} $p(1 - p) \times 0.3198 = 0.685 \times 0.315 \times 0.3198 = 0.069$: a
dalen more on the bus fare raises her probability of driving by about 7 points.

\textbf{8.} The marginal effect gives $0.2158 \times (-0.0537) \times 10 =
-0.116$. Exactly: $x'\hat\beta$ falls by 0.537 to 0.242, and $\Lambda(0.242) = 0.560$,
a fall of 0.125. The probability curve is not straight; the marginal effect is
its slope at the starting point only.

\textbf{9.} $60 \times 0.0537/0.3198 = 10.1$ dalens an hour.

\textbf{10.} $60 \times 0.0730/0.3335 = 13.1$ dalens an hour.
% figdata: commuters_mnl

\textbf{11.} Both are estimates with sampling error, and the binary model uses
fewer commuters and only the car--bus comparisons; its cost and time
coefficients have larger standard errors (0.0546 and 0.0114 against 0.0467
and 0.0073).

\textbf{12.} $20{,}000 \times 5 \times 220 = 22{,}000{,}000$ minutes, or
$366{,}667$ hours, times 13.13 dalens: about \dal{4.8} million a year.

\textbf{13.} $\Pr(\text{bicycle}) = \eu^{V_b}/(\eu^{V_c} + \eu^{V_u} + \eu^{V_b})$,
with $V_c = 0.4053 - 0.3335\,\text{cost}_c - 0.0730\,\text{time}_c$, $V_u =
-0.3335 \times 2.5 - 0.0730\,\text{time}_u$ and $V_b = -0.6719 -
0.0730\,\text{time}_b$.

\textbf{14.} 0.448, 0.213 and 0.339 in both columns. % figdata: commuters_shares
With a constant for every mode but the base, the score equations for the
constants say that fitted and observed shares are equal.

\textbf{15.} At equal cost and time, the bicycle is less attractive than the bus
by 0.672 in utility: worth $0.672/0.3335 = 2.0$ dalens a trip, or $0.672/0.0730
= 9.2$ minutes. It stands for effort, weather and danger.

\textbf{16.} $\eu^{-2} = 0.1353$, $\eu^{-3} = 0.0498$, $\eu^{-2.5} = 0.0821$, sum
0.2672: car 0.506, bus 0.186, bicycle 0.307.

\textbf{17.} $\eu^{-2}/\eu^{-3} = \eu = 2.72$. Without the bicycle it is still
2.72 (IIA): car 0.731, bus 0.269.

\textbf{18.} Yes, by assumption. A minute on a bicycle is effort and exposure,
not a seat; a mode-specific time coefficient would let the data say whether it
costs more.

\textbf{19.} 0.297, almost three commuters in ten.

\textbf{20.} Car: $0.448 - 0.319 = 0.129$ points, 28.8~\% of its share. Bus:
$0.213 - 0.143 = 0.070$ points, 32.9~\%. Bicycle: $0.339 - 0.241 = 0.098$ points,
28.9~\%. Most riders come from the car; the largest percentage from the bus.

\textbf{21.} For each commuter, the express takes from every mode in proportion
to that commuter's probabilities (IIA). The percentages differ across modes only
because the express is more attractive to some commuters, those for whom the
bus is already a good option.

\textbf{22.} From the ordinary bus: the express is the same product, only
faster. Commuters who drive because they dislike buses will mostly keep
driving.

\textbf{23.} With the two buses in one nest, tastes for bus travel are shared
by the two options, and the express takes most of its riders from the ordinary
bus; the bus-type share grows less, and the car and bicycle lose fewer riders.

\textbf{24.} Survey the same commuters after the express opens, compare each
mode's forecast share with the realised one, and tabulate forecast against
actual choices.

\textbf{25.} The commuters value travel time at about 13.1 dalens an hour. The
three-way logit forecasts that an express line 10 minutes faster than the bus
would carry 29.7~\% of commuters; because of IIA the forecast takes too many from
the car and the bicycle and too few from the ordinary bus, and is probably too
high.
\end{solution}
